\documentclass[12pt]{article}
\usepackage[utf8]{inputenc}
\usepackage[T1]{fontenc}
\usepackage[margin=1in]{geometry}
\usepackage{setspace}
\usepackage{xcolor}
\usepackage[colorlinks=true, urlcolor=blue!60!black, linkcolor=blue!60!black]{hyperref}
\setlength{\parindent}{2em}
\setlength{\parskip}{1.25em}
\renewcommand{\baselinestretch}{1.0}
\begin{document}
\begin{center}
{\LARGE \textit{Penn Central Transportation Co.\ v.\ New York City}, 438 U.S.\ 104 (1978)\par}
\vspace{0.8em}
{\large Case Brief --- Four-Page Edition\par}
\end{center}
\vspace{1em}

\section{Memory Jogger}

Upholding New York City's refusal to allow an office tower above Grand Central Terminal, the Court held 6--3 that no taking occurred and set the default test for regulatory takings: an ad hoc weighing of economic impact, interference with investment-backed expectations, and the character of the government action, measured against the parcel as a whole (\href{https://supreme.justia.com/cases/federal/us/438/104/}{Justia}).

\section{Detailed Case Facts}

New York City's 1965 Landmarks Preservation Law let an 11-member commission designate ``landmarks,'' whose owners had to keep exteriors ``in good repair'' and obtain approval before exterior alterations. To soften the burden, zoning rules gave landmark owners Transferable Development Rights (TDRs) to move unused air rights to nearby lots. In 1967 the Commission designated Grand Central Terminal, a 1913 Beaux-Arts station owned by Penn Central, as a landmark; Penn Central did not seek review. In 1968 Penn Central leased the airspace to UGP Properties for an office tower, at \$1 million a year during construction and at least \$3 million a year afterward. The Commission rejected both of architect Marcel Breuer's zoning-compliant designs, a 55-story tower cantilevered above the facade and a 53-story tower requiring partial demolition, calling the first an ``aesthetic joke.''

\section{Procedural History}

Without seeking review of the denials or proposing a smaller building, Penn Central and UGP sued in state court, alleging a taking and a due process violation. The trial court granted relief; the Appellate Division reversed; and the New York Court of Appeals affirmed, finding that Penn Central could still earn a reasonable return and held valuable TDRs. The Supreme Court affirmed on June 26, 1978.

\section{Judicial Votes}

Affirmed, \textbf{6--3}.
\begin{itemize}
\item \textbf{Majority:} Justice Brennan, joined by Justices Stewart, White, Marshall, Blackmun, and Powell.
\item \textbf{Concurrences:} None.
\item \textbf{Dissent:} Justice Rehnquist, joined by Chief Justice Burger and Justice Stevens.
\end{itemize}

\section{Holding}

Applying the Landmarks Law to Grand Central Terminal was not a taking. Regulatory takings claims require an ad hoc, factual inquiry weighing (1) the regulation's economic impact on the claimant, (2) its interference with distinct investment-backed expectations, and (3) the character of the governmental action, measured against the ``parcel as a whole'' rather than severed air rights. Because the law preserved the Terminal's profitable existing use, allowed a reasonable return, did not bar a smaller harmonizing addition, and conferred valuable TDRs, no compensation was due.

\section{Analysis of Opinions}

\textbf{Majority (Brennan).} The Takings Clause bars forcing a few people to bear public burdens that ``in fairness and justice'' the public should bear (\textit{Armstrong}), but there is no ``set formula'' for when regulation ``goes too far'' (\textit{Mahon}). A taking is more readily found for a physical invasion (\textit{Causby}) than for a program adjusting economic benefits and burdens. Brennan rejected four theories: Penn Central could not sever its air rights; the law was a comprehensive plan covering over 400 landmarks and 31 districts, not ``reverse spot'' zoning; general-welfare laws like zoning (\textit{Euclid}) and nuisance rules (\textit{Hadacheck}) often burden some owners more than others; and the City did not appropriate the airspace for its own use. Applying the factors, the law left the Terminal's 65-year primary use intact with a ``reasonable return,'' Penn Central never proposed a smaller structure, and TDRs usable on at least eight nearby properties mitigated the burden.

\textbf{Dissent (Rehnquist).} Zoning creates ``an average reciprocity of advantage,'' but New York singled out less than 0.1\% of its buildings, imposing multimillion-dollar losses and an affirmative upkeep duty backed by criminal penalties. The tower was zoning-compliant and no nuisance; Penn Central was burdened because ``too good a job was done'' in building the Terminal. He would find a taking and remand on whether TDRs were the ``full and perfect equivalent'' of compensation.

\textbf{Later.} \textit{Loretto} (1982) and \textit{Lucas} (1992) created \textit{per se} rules for permanent physical occupations and total wipeouts; \textit{Murr v.\ Wisconsin} (2017) refined the parcel; and \textit{Cedar Point} (2021) and \textit{Tyler} (2023) expanded physical and appropriative takings.

\section{Examples: Future Applications}

\begin{enumerate}
\item \textbf{Same-Side: A Tower Atop a Landmark Warehouse.} Denying a 30-story tower on a fully leased landmark earning 7\%, with TDRs, is no taking; the owner cannot sever air rights.
\item \textbf{Same-Side: A Wetlands Buffer.} Barring construction on the rear acre of a two-acre home lot is a partial burden on the parcel as a whole.
\item \textbf{Opposite-Side: A Public Easement and Arbitrary Fee.} A public sightseeing easement on a private dock is a \textit{per se} taking under \textit{Cedar Point}, and a flat fee unrelated to a permit's impact fails \textit{Nollan}/\textit{Dolan} after \textit{Sheetz} (2024).
\item \textbf{Opposite-Side: A Total Building Ban.} A setback law leaving beachfront lots with no economically beneficial use is a categorical taking under \textit{Lucas}.
\item \textbf{Fence-Sitter: Landmarking One Small Building Without TDRs.} A 75 percent loss soon after purchase, with no TDRs, pits frustrated expectations against a still-viable use.
\end{enumerate}

\section{Critique}

\textbf{Indeterminacy.} Critics note that \textit{Penn Central} assigns no weights to its factors, making outcomes unpredictable and favoring regulators whenever any viable use remains.

\textbf{The denominator and TDRs.} Defining the ``parcel as a whole'' is often outcome-determinative, yet \textit{Penn Central} and \textit{Murr} leave it fluid; counting TDRs risks letting governments offset takings with speculative credits instead of cash.

\section{Key Quotations}

\begin{quote}
``In engaging in these essentially ad hoc, factual inquiries, the Court's decisions have identified several factors that have particular significance. The economic impact of the regulation on the claimant and, particularly, the extent to which the regulation has interfered with distinct investment-backed expectations are, of course, relevant considerations. \dots\ So, too, is the character of the governmental action. A `taking' may more readily be found when the interference with property can be characterized as a physical invasion by government \dots\ than when interference arises from some public program adjusting the benefits and burdens of economic life to promote the common good.''
\end{quote}
\noindent --- Justice Brennan, for the Court, 438 U.S.\ at 124.

\begin{quote}
```Taking' jurisprudence does not divide a single parcel into discrete segments and attempt to determine whether rights in a particular segment have been entirely abrogated. In deciding whether a particular governmental action has effected a taking, this Court focuses rather both on the character of the action and on the nature and extent of the interference with rights in the parcel as a whole---here, the city tax block designated as the `landmark site.'\,''
\end{quote}
\noindent --- Justice Brennan, for the Court, 438 U.S.\ at 130--31.

\begin{quote}
``Appellees have imposed a substantial cost on less than one one-tenth of one percent of the buildings in New York City for the general benefit of all its people. It is exactly this imposition of general costs on a few individuals at which the `taking' protection is directed.''
\end{quote}
\noindent --- Justice Rehnquist, dissenting, 438 U.S.\ at 147.

\end{document}
